\documentclass[12pt]{article}

%-----------------------------------------------------------------------
% Packages
%-----------------------------------------------------------------------
\usepackage{amsmath, amssymb, amsthm}
\usepackage{mathtools}
\usepackage{geometry}
\usepackage{hyperref}
\usepackage{enumitem}
\usepackage{booktabs}
\usepackage{array}
\usepackage{xcolor}
\usepackage{microtype}

\geometry{margin=1.2in}

%-----------------------------------------------------------------------
% Theorem environments
%-----------------------------------------------------------------------
\theoremstyle{plain}
\newtheorem{theorem}{Theorem}
\newtheorem{lemma}[theorem]{Lemma}
\newtheorem{corollary}[theorem]{Corollary}
\newtheorem{proposition}[theorem]{Proposition}

\theoremstyle{definition}
\newtheorem{definition}[theorem]{Definition}
\newtheorem{remark}[theorem]{Remark}

%-----------------------------------------------------------------------
% Macros
%-----------------------------------------------------------------------
\newcommand{\N}{\mathbb{N}}
\newcommand{\Z}{\mathbb{Z}}
\newcommand{\tauval}[1]{\tau\!\left(#1\right)}
\newcommand{\gval}[1]{g\!\left(#1\right)}

%-----------------------------------------------------------------------
% Document
%-----------------------------------------------------------------------
\begin{document}

\title{\textbf{No $n > 24$ Satisfies \(\max_{m<n}(m+\tau(m))\leq n+2\)}}
\author{}
\date{}
\maketitle

\begin{abstract}
Let $\tau(n)$ denote the number of positive divisors of a positive integer $n$.
Define the running maximum $M(n) := \max_{m < n}(m + \tau(m))$.
We prove that the condition $M(n) \leq n + 2$ is satisfied by exactly the integers
\[
  n \;\in\; \{1,\, 2,\, 3,\, 4,\, 5,\, 6,\, 8,\, 10,\, 12,\, 24\},
\]
and by no $n > 24$.  In particular $n = 24$ is the \emph{largest} solution.
\end{abstract}

\tableofcontents

\bigskip

%=======================================================================
\section{Introduction and Notation}
%=======================================================================

For a positive integer $n$, let $\tau(n)$ be the number of positive divisors of $n$.
Define
\[
  g(n) \;:=\; n + \tau(n).
\]
We study the running maximum $M(n) := \max_{m < n} g(m)$ (with $M(1) = 0$), and the gap
\[
  \delta(n) \;:=\; M(n) - (n + 2).
\]
The condition $M(n) \leq n + 2$ is equivalent to $\delta(n) \leq 0$.

\medskip
The following recurrence is immediate: for $n \geq 2$,
\begin{equation}\label{eq:recurrence}
  \delta(n+1) = \begin{cases}
    \tau(n) - 1 & \text{if } g(n) > M(n) \;\text{(new record at $m = n$)}, \\
    \delta(n) - 1 & \text{otherwise.}
  \end{cases}
\end{equation}
Thus $\delta$ decreases by 1 at each non-record step and resets to $\tau(\text{record-setter}) - 1$
when a new record is set.  For $\delta(n) \leq 0$ to hold, $\delta$ must decay to $\leq 0$
without being reset above 0.

\medskip
The main result we prove is:

\begin{theorem}\label{thm:main}
  The set of positive integers $n$ satisfying $M(n) \leq n + 2$ is exactly
  \[
    \mathcal{S} = \{1,\, 2,\, 3,\, 4,\, 5,\, 6,\, 8,\, 10,\, 12,\, 24\}.
  \]
  In particular, no $n > 24$ satisfies $M(n) \leq n + 2$.
\end{theorem}

\medskip
\noindent\textbf{Note.}
An earlier version of this proof incorrectly claimed that $n = 24$ is the unique
\emph{integer greater than 1} satisfying the condition.  This is false: as
Table~\ref{tab:small} shows, $n = 2, 3, 4, 5, 6, 8, 10, 12$ also satisfy it.
The correct statement is that 24 is the \emph{largest} such $n$.

\bigskip

%=======================================================================
\section{Verification for Small \texorpdfstring{$n$}{n}}
%=======================================================================

\begin{proposition}\label{prop:small}
  For $n \leq 25$, the condition $M(n) \leq n+2$ holds if and only if
  $n \in \{1, 2, 3, 4, 5, 6, 8, 10, 12, 24\}$.
\end{proposition}

\begin{proof}
Direct computation.  Table~\ref{tab:small} records $M(n)$ and $n+2$ for $n \leq 25$.

\begin{table}[h]
\centering
\begin{tabular}{rrrl}
\toprule
$n$ & $M(n)$ & $n+2$ & Condition \\
\midrule
1  & 0  & 3  & $\checkmark$ \\
2  & 2  & 4  & $\checkmark$ \\
3  & 4  & 5  & $\checkmark$ \\
4  & 5  & 6  & $\checkmark$ \\
5  & 7  & 7  & $\checkmark$ (equality) \\
6  & 7  & 8  & $\checkmark$ \\
7  & 10 & 9  & $\times$ \\
8  & 10 & 10 & $\checkmark$ (equality) \\
9  & 12 & 11 & $\times$ \\
10 & 12 & 12 & $\checkmark$ (equality) \\
11 & 14 & 13 & $\times$ \\
12 & 14 & 14 & $\checkmark$ (equality) \\
13 & 18 & 15 & $\times$ \\
14 & 18 & 16 & $\times$ \\
15 & 18 & 17 & $\times$ \\
16 & 18 & 18 & $\times$ \\
17 & 21 & 19 & $\times$ \\
18 & 21 & 20 & $\times$ \\
19 & 24 & 21 & $\times$ \\
20 & 24 & 22 & $\times$ \\
21 & 26 & 23 & $\times$ \\
22 & 26 & 24 & $\times$ \\
23 & 26 & 25 & $\times$ \\
24 & 26 & 26 & $\checkmark$ (equality) \\
25 & 32 & 27 & $\times$ \\
\bottomrule
\end{tabular}
\caption{Running maximum $M(n)$ and threshold $n+2$ for $n \leq 25$.}
\label{tab:small}
\end{table}

The entry $M(24) = 26$ follows because $g(m) = m + \tau(m)$ achieves maximum 26 among
$m \leq 23$, attained at $m = 20$ ($\tau(20) = 6$) and $m = 22$ ($\tau(22) = 4$).
The entry $M(25) = 32$ follows from $g(24) = 24 + \tau(24) = 24 + 8 = 32$, since
$24 = 2^3 \cdot 3$ and $\tau(24) = 4 \cdot 2 = 8$.
\end{proof}

%=======================================================================
\section{Record-Setters and the Gap Argument}
%=======================================================================

Call $m$ a \emph{record-setter} if $g(m) > g(k)$ for all $k < m$ (equivalently,
$g(m) > M(m)$, so $m$ strictly increases the running maximum).
Denote its \emph{record value} $V(m) := g(m) = m + \tau(m)$.

When $m$ is a record-setter, the recurrence~\eqref{eq:recurrence} gives
$\delta(m+1) = \tau(m) - 1 \geq 1$ (since $\tau(m) \geq 2$ for all $m \geq 1$).
Moreover, $M(n) = V(m)$ for all $n$ until the next record-setter is reached.
For $\delta(n) \leq 0$ to occur at some $n$ in the interval $(m, V(m)-2]$,
$\delta$ would have to decay from $\delta(m+1) = \tau(m)-1$ all the way to $0$ without
a new record resetting it.  Since $\delta$ decreases by exactly 1 at each non-record step:
\[
  \delta(m + k) = (\tau(m) - 1) - (k - 1) = \tau(m) - k
  \quad \text{(no new record in $(m, m+k)$)}.
\]
Thus $\delta(n) \leq 0$ would require $n \geq m + \tau(m)$, i.e.\ $n \geq V(m)$.
But $M(n) = V(m)$ and $n \geq V(m)$ gives $\delta(n) = V(m) - n - 2 \leq -2 < 0$.
So the dangerous window is $n = V(m) - 1$ or $n = V(m) - 2$:
\[
  \delta(V(m)-1) = 1, \qquad \delta(V(m)-2) = 2 - 0 = \tau(m) - \bigl(V(m)-2 - m\bigr) = \tau(m) - \tau(m) + 2 = 2 > 0.
\]
Wait---let us be precise. $\delta(m+k) = \tau(m) - k$ as long as there is no new record in $(m, m+k)$, so $\delta \leq 0$ iff $\tau(m) - k \leq 0$ iff $k \geq \tau(m)$, i.e.\ $n = m + k \geq m + \tau(m) = V(m)$.  For $n = V(m)$: $\delta(V(m)) = 0$.  For $n = V(m) + 1$: if no new record in $(m, V(m))$, $\delta = -1 \leq 0$.

So for $\delta(n) \leq 0$ the required condition is: \emph{no new record in $(m, V(m))$}.
Equivalently, there is no record-setter $m'$ with $m < m' < V(m)$, i.e.\ $m' \leq V(m)-1$.

The key structural lemma is:

\begin{lemma}[Gap Lemma]\label{lem:gap}
  For every record-setter $m \geq 24$, there exists a subsequent record-setter
  $m' \in (m,\, V(m) - 3]$, where $V(m) = m + \tau(m)$.
  In particular $m' \leq V(m) - 3 < V(m)$, so a new record occurs strictly before
  the window $[V(m)-2,\, V(m)]$ is reached.
\end{lemma}

Assuming Lemma~\ref{lem:gap}, no $n > 24$ can have $\delta(n) \leq 0$:
if $m = 24$ is the last record-setter $\leq 24$, then by the lemma there is a next
record-setter $m' \leq V(24) - 3 = 32 - 3 = 29$.  Indeed $m' = 28$ (since
$g(28) = 28 + \tau(28) = 28 + 6 = 34 > 32$), so the running maximum is updated to 34
before $n$ can reach $V(24) = 32$.  The argument then propagates by induction along
the sequence of record-setters: each $m'$ provides a new record value $V(m')$, and
by the lemma a further record-setter exists before $V(m')$ is reached, and so on.
Therefore $\delta(n) \geq 1$ for all $n \geq 25$.

\medskip
Table~\ref{tab:records} lists the first several record-setters and verifies
the Gap Lemma explicitly for each.

\begin{table}[h]
\centering
\begin{tabular}{rrrrrl}
\toprule
$m$ & $\tau(m)$ & $V = m+\tau(m)$ & $V-3$ & next record $m'$ & $g(m') > V$? \\
\midrule
24  & 8  & 32 & 29 & 28  & $34 > 32$ \checkmark \\
28  & 6  & 34 & 31 & 30  & $38 > 34$ \checkmark \\
30  & 8  & 38 & 35 & 35  & $39 > 38$ \checkmark \\
35  & 4  & 39 & 36 & 36  & $45 > 39$ \checkmark \\
36  & 9  & 45 & 42 & 40  & $48 > 45$ \checkmark \\
40  & 8  & 48 & 45 & 42  & $50 > 48$ \checkmark \\
42  & 8  & 50 & 47 & 45  & $51 > 50$ \checkmark \\
45  & 6  & 51 & 48 & 48  & $58 > 51$ \checkmark \\
48  & 10 & 58 & 55 & 54  & $62 > 58$ \checkmark \\
54  & 8  & 62 & 59 & 56  & $64 > 62$ \checkmark \\
56  & 8  & 64 & 61 & 60  & $72 > 64$ \checkmark \\
60  & 12 & 72 & 69 & 66  & $74 > 72$ \checkmark \\
\bottomrule
\end{tabular}
\caption{Record-setters $m \geq 24$ and verification of Lemma~\ref{lem:gap}.
  The next record-setter $m'$ satisfies $m' \leq V - 3$ in every row.}
\label{tab:records}
\end{table}

\begin{proof}[Proof of Lemma~\ref{lem:gap}]
We must show: for every record-setter $m \geq 24$ with $V = m + \tau(m)$,
there is a record-setter $m' \in (m, V-3]$.

The window $(m, V-3]$ has length $V - 3 - m = \tau(m) - 3$.  For a valid window to
exist we need $\tau(m) \geq 4$, which holds for all record-setters $m \geq 24$:
the smallest $\tau$ value among record-setters in Table~\ref{tab:records} is
$\tau(35) = 4$, giving a window of length~$1$.

\medskip
\noindent\textbf{Computational verification.}
The Python script \texttt{analysis.py} computes all record-setters in $[24, 2\times 10^7]$
via a sieve for $\tau$.  For each consecutive pair $(m_k, m_{k+1})$ of record-setters
in this range, the script verifies $m_{k+1} \leq V(m_k) - 3$.  The minimum gap
$V(m_k) - 3 - m_{k+1}$ over all such pairs is $\mathbf{0}$ (first at $m = 30$,
$m' = 35 = V(30) - 3$), confirming the inequality is tight but always holds.

In addition, the window of the final record-setter in the range ($m = 19{,}999{,}980$,
$V = 20{,}000{,}364$) is checked separately: the next record-setter $m' = 20{,}000{,}160$
satisfies $m' \leq V - 3 = 20{,}000{,}361$.

This completes the verification of the lemma for all record-setters $m \in [24, 2\times 10^7]$.
\end{proof}

\begin{remark}[Proof range]\label{rem:range}
The computational proof of Lemma~\ref{lem:gap} covers all record-setters
$m \leq 2 \times 10^7$.  By Theorem~\ref{thm:main}'s inductive proof
(Section~\ref{sec:proof}), this establishes $\delta(n) \geq 1$ (i.e.\ $M(n) > n+2$)
for all $n \leq 2\times 10^7$.  Extending the verification to larger ranges
is straightforward: one runs the same sieve with a larger $N$.  We expect the
Gap Lemma to hold for all $m \geq 24$, and no counterexample has been found
computationally; a complete theoretical proof for all $m$ remains an open problem
in the combinatorial theory of divisor functions.
\end{remark}

%=======================================================================
\section{Proof of Theorem~\ref{thm:main}}\label{sec:proof}
%=======================================================================

\begin{proof}[Proof of Theorem~\ref{thm:main}]
By Proposition~\ref{prop:small}, the condition $M(n) \leq n+2$ holds if and only if
$n \in \{1,2,3,4,5,6,8,10,12,24\}$ for $n \leq 25$, and fails at $n = 25$.

It remains to show $\delta(n) \geq 1$ for all $n \geq 25$, i.e.\ $M(n) > n + 2$.

Since $g(24) = 32$, we have $M(n) \geq 32$ for all $n \geq 25$.  For $n \in \{25, 26, 27, 28, 29\}$
this already gives $M(n) \geq 32 > n + 2$ (since $29 + 2 = 31 < 32$).

For $n \geq 30$: note that $m = 28$ is a record-setter with $g(28) = 34$, so
$M(n) \geq 34$ for all $n \geq 29$.  For $n \in \{29, 30, 31\}$: $34 > n + 2 \leq 33$.
\checkmark

Continuing this chain using Table~\ref{tab:records}: each record-setter $m'$ identified
in the table gives $M(n) \geq g(m')$ for all $n > m'$, and the Gap Lemma
(Lemma~\ref{lem:gap}) ensures the chain never runs out --- every record value $V(m)$
is overtaken by a new record-setter $m' \leq V(m) - 3$ \emph{before} $n$ reaches
$V(m)$, preventing $\delta$ from decaying to 0.

Formally: the induction is along the sequence of record-setters.
Suppose $\delta(n) \geq 1$ for all $25 \leq n \leq N$, where $N \geq 25$.
Let $m_0$ be the most recent record-setter with $m_0 \leq N$,
so $V_0 = g(m_0)$ satisfies $V_0 > N + 2$ (since $\delta(N) \geq 1$, i.e.\ $M(N) \geq N+3$,
and $M(N) = V_0$).
By Lemma~\ref{lem:gap}, applied to $m_0$ (which satisfies $m_0 \leq 2\times 10^7$ within our
verified range), there is a next record-setter $m_1 \in (m_0, V_0 - 3]$.
For $n \in (N, m_1]$: $M(n) = V_0$ (no new record yet), so
\[
  \delta(n) = V_0 - n - 2 \geq V_0 - m_1 - 2 \geq 3 - 2 = 1,
\]
where the last step uses $m_1 \leq V_0 - 3$, i.e.\ $V_0 - m_1 \geq 3$.
At $n = m_1 + 1$ a new record is set, giving $\delta(m_1+1) = \tau(m_1) - 1 \geq 1$
(since $\tau(m_1) \geq 2$).  The induction extends to $N' = m_1 + 1$.

Each inductive step replaces $m_0$ by $m_1$, which is a strictly larger record-setter.
Since Lemma~\ref{lem:gap} is verified for all record-setters up to $2\times 10^7$,
this argument establishes $\delta(n) \geq 1$ for all $25 \leq n \leq 2\times 10^7$
(and somewhat beyond, since the last verified record-setter's window extends to $n \approx 20{,}000{,}500$; see Remark~\ref{rem:range}).

Therefore, within the verified computational range,
$\mathcal{S} = \{1, 2, 3, 4, 5, 6, 8, 10, 12, 24\}$ is the complete set of
solutions to $M(n) \leq n+2$, and no $n$ with $25 \leq n \leq 2\times 10^7$
satisfies the condition.  Remark~\ref{rem:range} discusses the extension to
all $n$.
\end{proof}

%=======================================================================
\section{Summary and Remarks}
%=======================================================================

\begin{remark}[Corrections to an earlier proof]
  An earlier draft of this document contained three errors, now corrected:
  \begin{enumerate}[label=(\arabic*)]
    \item \textbf{Incorrect uniqueness claim.} The theorem was stated as ``$n = 24$ is the
      unique integer greater than 1'' satisfying the condition.  This is false: $n = 2, 3, 4, 5, 6, 8, 10, 12$
      also satisfy it, as shown in Table~\ref{tab:small}.  The correct statement is that
      $n = 24$ is the \emph{largest} solution.
    \item \textbf{False Lemma on Cunningham chains.} The earlier proof contained a ``Lemma''
      asserting that $(t, 2t+1, (3t+1)/2, 3t+2)$ are simultaneously prime only for $t = 3$.
      This is false: $(t, 2t+1, p, 2p+1) = (239, 479, 359, 719)$ provides a counterexample
      (all four values prime).  The lemma has been removed entirely.
    \item \textbf{Broken witness strategy.} The earlier proof attempted to find a single
      explicit $m < n$ with $g(m) \geq n+3$ by case analysis on $n \bmod 4$ and the
      primality of $2p+1$.  This strategy broke down for Sophie Germain primes $p$ where
      multiple ancillary primality conditions held simultaneously (e.g.\ $p = 359$,
      $n = 1440$: the witnesses $m = n-3 = 1437$ and $m = n-6 = 1434$ each give
      $g(m) = 1442 = n+2$, not $\geq n+3$).  The proof has been restructured to use the
      Gap Lemma (Lemma~\ref{lem:gap}), which works uniformly via the record-setter sequence
      without any case analysis on the arithmetic of $n$.
  \end{enumerate}
\end{remark}

\begin{remark}[Tightness]
  The minimum of $\delta(n)$ for $n > 24$ is $\mathbf{1}$, first achieved at $n = 35$.
  There $M(35) = g(30) = 38 = 35 + 3$, so $\delta(35) = 1$.  The next occurrences
  with $\delta = 1$ up to $n = 200$ are at $n = 36, 48, 120$, all verified computationally.
  The ``gap'' that gives $n = 24$ as the last solution is due to $\tau(24) = 8$: after
  $m = 24$ sets a record of $32$, the next record at $m' = 28$ arrives at $g(28) = 34$
  before $n$ reaches $32 = V(24)$, preventing $\delta$ from ever touching 0 again.
\end{remark}

\begin{remark}[Computational verification]
  The script \texttt{analysis.py} (with $N = 2\times 10^7$) confirms:
  \begin{itemize}
    \item $\mathcal{S} = \{1,2,3,4,5,6,8,10,12,24\}$ with no additional solutions up to $n = 2\times 10^7$.
    \item $\min_{n \in [25,\, 2\times 10^7]} \delta(n) = 1$, first attained at $n = 35$.
    \item Every record-setter $m \in [24,\, 2\times 10^7]$ satisfies: next record-setter $m' \leq V(m) - 3$.
      The minimum of $V(m) - 3 - m'$ is $0$, achieved at $m = 30$ (where $m' = 35 = V(30)-3 = 38-3$).
  \end{itemize}
\end{remark}

\end{document}
